% INTRODUCTION - Auto-updated section
% Last generated: 2026-01-15T05:34:34.872Z

\subsection{Motivation}

The emergence of Decentralized Autonomous Organizations (DAOs) represents one of the most significant experiments in organizational design since the joint-stock corporation. By encoding governance rules in smart contracts and distributing decision-making authority among token holders, DAOs promise to enable new forms of coordination at scale without traditional hierarchical structures \citep{buterin2014dao, wright2015decentralized}.

However, the design space for DAO governance mechanisms is vast and poorly understood. Practitioners face fundamental questions: What quorum threshold balances legitimacy with operational efficiency? How should voting power be allocated to prevent plutocratic capture while maintaining stakeholder alignment? When do delegation systems improve or harm collective decision-making?

Traditional approaches to answering these questions---theoretical analysis and empirical observation of deployed systems---each have limitations. Game-theoretic analysis often requires simplifying assumptions that may not hold in practice. Empirical studies of live DAOs conflate mechanism effects with confounding factors like market conditions, community culture, and historical path dependencies.

\subsection{Contributions}

We present a multi-agent simulation framework that enables systematic, controlled experimentation with DAO governance mechanisms. Our contributions include:

\begin{enumerate}
    \item \textbf{Simulation Framework}: An open-source, extensible simulator supporting multiple governance mechanisms (token voting, quadratic voting, conviction voting), agent heterogeneity, delegation, and realistic behavioral models.

    \item \textbf{Experimental Infrastructure}: A research CLI enabling reproducible batch experiments with parameter sweeps, checkpoint/resume capabilities, and export to standard data analysis tools.

    \item \textbf{Empirical Analyses}: A reproducible experiment suite and analysis pipeline; results are reported as runs complete (21402 total runs across 10 experiment sets in this version).

    \item \textbf{Theoretical Integration}: A framework connecting multi-agent systems, mechanism design, and computational social choice theory to DAO governance analysis.
\end{enumerate}

\subsection{Research Questions}

This paper addresses the following research questions:

\begin{description}
    \item[RQ1:] How do quorum thresholds affect proposal passage rates and governance efficiency?
    \item[RQ2:] What is the relationship between DAO size and voter participation?
    \item[RQ3:] How do different voting mechanisms (token, quadratic, conviction) affect wealth concentration and minority representation?
    \item[RQ4:] Under what conditions do delegation systems improve collective decision-making?
\end{description}

\subsection{Paper Organization}

Section~\ref{sec:background} reviews related work in multi-agent systems, mechanism design, and DAO governance. Section~\ref{sec:theory} develops our theoretical framework. Section~\ref{sec:architecture} describes the simulation architecture. Section~\ref{sec:methodology} details our experimental methodology. Section~\ref{sec:results} presents findings from 0 experimental configurations. Section~\ref{sec:discussion} interprets results and discusses implications. Section~\ref{sec:limitations} acknowledges limitations and outlines future work. Section~\ref{sec:conclusion} concludes.
